% Fuente conservada íntegra: originales/cuaderno/NOTA.md, §§1--7.
\section{Memoria cronológica, balance y respuesta de vacancias}
\label{md:cron:seccion}

La construcción siguiente recibe el factor simbólico de capacidad y estudia
sus lectores posteriores. La razón interna $3^6/10^3=729/1000$ y su
calendario pertenecen a la construcción precedente
(secciones~\ref{mono-ampl:relojes} y~\ref{vac-capacidad-trit}).
La arquitectura APP--TRIT--TPK, las regiones, el retorno nonádico y el
registro $K$ conservan su posición causal anterior. Los logaritmos empleados
en el lector de capacidad describen dicho factor y no seleccionan de nuevo
las regiones de $\pi,\varphi,e,\alpha$.
Las proposiciones son consecuencias focales sobre observación y memoria;
su alcance se distingue del estado TPK completo y de una atribución de
prioridad histórica a resultados generales sobre sucesiones mecánicas o
información.

\subsection{Objeto, medida y dos lectores}
\label{md:cron:lectores}

Sea $\delta=\log_{10}(10/9)$ la pendiente del calendario de capacidad,
y sea $\theta$ una fase uniformemente distribuida en $[0,1)$.
La medida representa incertidumbre sobre el origen de lectura, mientras el
mecanismo permanece determinado. Definimos
\begin{equation}
 \begin{aligned}
 z_j(\theta)&=\lfloor(j+1)\delta+\theta\rfloor
             -\lfloor j\delta+\theta\rfloor,\\
 W_n(\theta)&=(z_0(\theta),\ldots,z_{n-1}(\theta)),\\
 B_n(\theta)&=\sum_{j=0}^{n-1}z_j(\theta).
 \end{aligned}
 \label{md:cron:dos-lectores}
\end{equation}
El bloque $W_n$ conserva la cronología y $B_n$ conserva el balance.
Para una fase exacta y una regla exacta conocidas, ambos objetos son
deterministas; las entropías se refieren expresamente al conjunto de fases
posibles.

La telescopía da
\begin{equation}
 B_n(\theta)=\lfloor n\delta+\theta\rfloor.
 \label{md:cron:telescopia}
\end{equation}
Por tanto, $B_n$ admite sólo los valores $\lfloor n\delta\rfloor$ y
$\lceil n\delta\rceil$. Los puntos $\{-j\delta\}$, con $j=0,\ldots,n$,
dividen el círculo en $n+1$ intervalos de fase. Cada intervalo porta una
secuencia distinta de longitud $n$. En efecto, los prefijos acumulados
$\lfloor k\delta+\theta\rfloor$, para $k=1,\ldots,n$, son monótonos en
$\theta$, y cada frontera cambia uno de ellos; la secuencia reconstruye
todos esos prefijos.

Los extremos tienen medida cero. Una convención semiabierta resuelve su
asignación y conserva las entropías.

\subsection{Información creciente y tasa entrópica nula}
\label{md:cron:crecimiento}

Sean $\ell_0,\ldots,\ell_n$ las longitudes de los intervalos de fase y
\begin{equation}
 H_n=-\sum_{j=0}^{n}\ell_j\log_2\ell_j.
 \label{md:cron:entropia-bloque}
\end{equation}

\begin{proposition}
\label{md:cron:proposicion-tasa}
Se cumplen simultáneamente
\begin{equation}
 H_n\longrightarrow+\infty,\qquad
 \frac{H_n}{n}\longrightarrow0.
 \label{md:cron:dos-limites}
\end{equation}
\end{proposition}
\begin{proof}
La cota superior es $H_n\leq\log_2(n+1)$. Si
$\ell_{\max}(n)=\max_j\ell_j$, entonces
$H_n\geq-\log_2\ell_{\max}(n)$.
La pendiente $\delta$ es irracional: si $\delta=p/q$ con $q>0$,
se tendría $(10/9)^q=10^p$, lo cual contradice la valuación del primo $3$.
Por irracionalidad, los puntos de la órbita son densos. Las particiones
sucesivas tienen malla máxima tendente a cero: para cada $\epsilon>0$
existe un tramo finito $\epsilon$-denso, y los tramos posteriores sólo
subdividen sus intervalos. Por tanto, la cota inferior diverge. La cota
superior dividida por $n$ tiende a cero.
\end{proof}

Las longitudes de los intervalos no son iguales en general. La prueba
establece las cotas anteriores, sin afirmar
$H_n=\log_2(n+1)$ ni una ley asintótica logarítmica exacta que requiriese
condiciones diofánticas adicionales.
La entropía topológica del factor también es cero porque
$\log(n+1)/n\to0$. La información del bloque y su tasa son cantidades
diferentes. El crecimiento del archivo de una trayectoria conocida no
equivale, por sí solo, a producción de información estadística independiente.

\subsection{Información cronológica no determinada por el balance}
\label{md:cron:condicional}

\begin{corollary}
\label{md:cron:corolario-condicional}
Para el mismo conjunto de fases,
\begin{equation}
 H(W_n\mid B_n)=H_n-H(B_n)
 \geq H_n-1\longrightarrow+\infty.
 \label{md:cron:entropia-condicional}
\end{equation}
\end{corollary}
\begin{proof}
$B_n$ es función de $W_n$, de modo que la regla de la cadena da la igualdad.
Su soporte tiene dos elementos, por lo que $H(B_n)\leq1$ bit.
Se aplica la proposición~\ref{md:cron:proposicion-tasa}.
\end{proof}

La entropía condicional que diverge es la media sobre los balances;
no se atribuye la misma incertidumbre a cada balance concreto.
Tampoco se atribuye esa pérdida al estado enriquecido completo: cuando la
memoria retenida recupera $W_n$, se tiene
\begin{equation}
 H(W_n\mid B_n,M_n)=0.
 \label{md:cron:restitucion-memoria}
\end{equation}
Por cardinalidad, al menos una de las dos fibras de $B_n$ contiene
$\lceil(n+1)/2\rceil$ secuencias o más. Esta cota combinatoria se distingue
de la cota entrópica.
El resultado precisa la relación entre valor y genealogía en este lector:
los dos balances posibles no agotan las historias temporales compatibles.

\subsection{Corriente centrada: balance acotado y actividad persistente}
\label{md:cron:corriente}

El lector centrado se expresa por
\begin{equation}
 J_j=z_j-\delta,\qquad
 \sum_{j=0}^{n-1}J_j=B_n-n\delta.
 \label{md:cron:corriente-definicion}
\end{equation}
En consecuencia,
$\left|\sum_{j=0}^{n-1}J_j\right|<1$ para toda ventana.
A la vez,
\begin{equation}
 \lim_{n\to\infty}\frac1n\sum_{j=0}^{n-1}J_j=0,\qquad
 \lim_{n\to\infty}\frac1n\sum_{j=0}^{n-1}J_j^2
 =\delta(1-\delta)>0.
 \label{md:cron:actividad}
\end{equation}
\begin{proof}
Se tiene $z_j^2=z_j$ y $B_n/n\to\delta$ por la cota de discrepancia.
La expansión
$(z_j-\delta)^2=z_j(1-2\delta)+\delta^2$, seguida de la sustitución del
promedio de $z_j$, da la segunda identidad.
El argumento vale para cualquier fase y no requiere un promedio aleatorio.
\end{proof}

El balance acumulado acotado, la actividad cuadrática positiva y la tasa de
entropía cero son, por tanto, compatibles. Las identidades conservan la
distinción entre el lector centrado y una representación de carga
elemental; tampoco establecen una producción energética sin balance.

Si una realización dimensional ya establecida toma
$I_j=(u_Q/t_0)J_j$, se obtiene
\begin{equation}
 I_{\mathrm{rms}}^2
 =\left(\frac{u_Q}{t_0}\right)^2\delta(1-\delta).
 \label{md:cron:eficaz}
\end{equation}
Si cada $I_j$ se mantiene durante un intervalo de duración $t_0$, este
promedio discreto coincide con el promedio temporal.
Como contraste físico posterior, una carga resistiva ideal
$R_{\mathrm{res}}>0$ tendría potencia media
\begin{equation}
 \overline P_{\mathrm{res}}
 =R_{\mathrm{res}}\left(\frac{u_Q}{t_0}\right)^2
 \delta(1-\delta)>0.
 \label{md:cron:potencia-resistiva}
\end{equation}
Es una predicción del modelo compuesto con esa carga. La resistencia y el
dispositivo conservan su especificación propia; no se derivan de la sola
sucesión de capacidad.

\subsection{Espectro de la corriente y orientación}
\label{md:cron:espectro}

La función de observación es
$f(\theta)=\mathbf1_{[1-\delta,1)}(\theta)-\delta$, con
$J_j=f(\theta+j\delta)$. Sus coeficientes de Fourier, para $m\neq0$, son
\begin{equation}
 a_m=\frac{\exp(2\pi i m\delta)-1}{2\pi i m},\qquad
 |a_m|^2=\frac{\sin^2(\pi m\delta)}{\pi^2m^2},\qquad
 a_0=0.
 \label{md:cron:fourier}
\end{equation}
Se obtienen integrando la indicatriz sobre su intervalo. La identidad de
Parseval da
\begin{equation}
 \sum_{m\neq0}|a_m|^2=\delta(1-\delta),
 \label{md:cron:parseval}
\end{equation}
en concordancia con la actividad cuadrática. Las frecuencias temporales
son $m\delta$ módulo $1$. La correlación estacionaria, con fase uniforme,
queda determinada por esas intensidades; éstas no contienen la fase inicial.
Una traslación multiplica $a_m$ por una fase unitaria, y la inversión
temporal intercambia las frecuencias opuestas.

La potencia espectral aislada no recupera la orientación ni el origen
temporal. Una lectura sensible a fase o una referencia marcada aporta
la información complementaria. La deducción conserva la distinción entre
esa referencia general y el registro $K$, así como entre ritmos simbólicos
y una realización experimental determinada.

\subsubsection{Referencia cíclica y transportes de igual intensidad espectral}
\label{md:cron:k-ciclico}

La propiedad cíclica del registro $K$ permite una consecuencia finita
independiente del espectro anterior. En $V_K=\mathbb R^{12}$, sea $S$ el
desplazamiento $(Sx)_i=x_{i+1}$ y
\begin{equation}
 O_K=(K,SK,\ldots,S^{11}K).
 \label{md:cron:matriz-ciclica}
\end{equation}
Se utiliza la invertibilidad de $O_K$ establecida por la referencia cíclica
del registro. Para cualquier operador real $H$ que conmute con $S$,
pongamos $y_+=HK$ e $y_-=H^TK$.

\begin{proposition}
\label{md:cron:proposicion-intensidades}
Las doce intensidades de Fourier de $y_+$ e $y_-$ coinciden.
Sin embargo, $y_+=y_-$ si y sólo si $H=H^T$.
Cada respuesta vectorial completa, junto con $K$ y la orientación de $S$,
recupera su operador.
\end{proposition}
\begin{proof}
$H$ es circulante. Sus multiplicadores de Fourier son $\lambda_m$ y los
de $H^T$ son sus conjugados, de modo que las dos respuestas tienen
intensidades $|\lambda_m|^2|\widehat K_m|^2$.
Si sus vectores coinciden, $(H-H^T)K=0$. La conmutación implica que
$H-H^T$ anula también $S^jK$ para todo $j$. Esos vectores forman una base;
por tanto, $H-H^T=0$. La implicación inversa es inmediata.
La recuperación se expresa por
\begin{equation}
 H=O_{y_+}O_K^{-1},\qquad
 H^T=O_{y_-}O_K^{-1}.
 \label{md:cron:restitucion-operador}
\end{equation}
\end{proof}

En particular, los desplazamientos $S$ y $S^{-1}$ producen respuestas con
el mismo espectro de potencia, pero distintas. Las autocorrelaciones exactas
del registro dan
\begin{equation}
 \|SK-S^{-1}K\|^2
 =2\bigl(\|K\|^2-\langle K,S^2K\rangle\bigr)
 =2175744>0.
 \label{md:cron:separacion-desplazamientos}
\end{equation}
La inversión racional de $O_K$ recupera de ambas respuestas los
coeficientes $e_1$ y $e_{11}$, respectivamente. $K$ se utiliza después
de su generación como referencia marcada. La intensidad espectral o el
balance escalar no sustituyen la respuesta vectorial.
El corolario describe el portador cíclico de doce componentes;
conserva la diferencia entre $S$ y la dinámica enriquecida completa,
así como los operadores adicionales de una conjugación física de carga.
Para un $H$ general, la transposición tampoco significa inversión temporal;
en el caso $S$, sí coincide con $S^{-1}$.

\subsection{Enlace operatorio electrónico}
\label{md:cron:electron}

La incidencia entre las fibras de suma y producto de APP sobre $D_9^3$
define $C$. Su vector trítico
$m_{\mathrm{ph}}=(1,-1,0)^3$ satisface
\begin{equation}
 Cm_{\mathrm{ph}}=C^Tm_{\mathrm{ph}}
 =-\frac12m_{\mathrm{ph}}.
 \label{md:cron:modo-tritico}
\end{equation}
El calendario inducido recorre los regímenes del mismo selector trítico;
la incidencia selecciona el plano principal de singularidad $1/2$.
En él,
\begin{equation}
 P=\begin{pmatrix}1&0\\0&0\end{pmatrix},\qquad
 Q=\begin{pmatrix}1/4&\sqrt3/4\\\sqrt3/4&3/4\end{pmatrix},
 \qquad J=\frac4{\sqrt3}[P,Q],\qquad J^2=-I.
 \label{md:cron:plano-electronico}
\end{equation}
La incidencia y la norma del conmutador se desarrollan en la
sección~\ref{vi:dep:prefactor}. Se recupera así una relación efectiva entre
organización temporal y bloque electrónico, conservando el selector,
sus etiquetas y la incidencia. La representación de carga y la realización
dimensional mantienen sus operadores propios: el enlace no asigna
directamente las cargas elementales opuestas a $z_j-\delta$.

\subsection{Información retenida y composición energética}
\label{md:cron:informacion-energia}

La recuperación de fibra establece
\begin{equation}
 H(X\mid q(X),M(X))=0
 \label{md:cron:inversa-enriquecida}
\end{equation}
cuando el lector enriquecido tiene inversa; una actualización biyectiva
conserva $H(X)$. Estas propiedades se reúnen en
\eqref{eq:ii-kb-aut-fibra} y~\eqref{eq:ii-kb-aut-reversible}.
La coexistencia de $H_n$ creciente y tasa cero es compatible con esas
propiedades y con las identidades de Shannon.

El trabajo de borrado depende de la información accesible retenida.
En el régimen ideal isotermo de memorias lógicas degeneradas y en el
límite asintótico apropiado, el coste por copia con información lateral es
\begin{equation}
 W_{\mathrm{borrado,ideal}}
 =k_BT\ln(2)\,H(X\mid M).
 \label{md:cron:borrado}
\end{equation}
Si $X$ se recupera desde $M$, la contribución lógica ideal puede ser cero.
El transporte físico, la lectura, el control y el ciclo completo conservan
balances propios; la igualdad $h_{\mathrm{top}}=0$ no los elimina.

La composición relevante conserva la secuencia
\[
 \begin{gathered}
 \text{cronología}\longrightarrow
 \text{registro de frontera}\longrightarrow
 \text{lector de corriente y estado}\\
 \longrightarrow
 \text{memoria retenida}\longrightarrow
 \text{balance térmico}.
 \end{gathered}
\]
El desarrollo anterior construye los dos primeros lectores y sus
consecuencias informacionales. Su alcance se distingue de la certificación
de un dispositivo físico completo y mantiene intacta la genealogía
constructiva de las constantes.
